\chapter{Implementations}

The repository carries two runnable implementations: a Python proof skeleton
that remains the authoritative oracle for namespace and logical-size parity,
and a native Rust read-only path that has been extended to FS-tree record-family
inventory. This chapter documents both and the discipline that keeps them
honest about what they prove.

\section{Python Proof Skeleton}

\subsection{Current Implementation Boundary}

The proof-backed parser skeleton is documented in
\path{docs/implementation/narrow-v1-proof-parser-skeleton.md}.

The execution path is:

\begin{enumerate}
  \item source gate
  \item scan-state summary
  \item proof raw-walk backend
  \item namespace entry normalization
  \item directory aggregate construction
  \item oracle diff
\end{enumerate}

\subsection{Regression Command}

From the repository root:

\begin{verbatim}
PYTHONPATH=src python3 -m apfs_fastindex.poc_regression
\end{verbatim}

The regression creates a temporary APFS image, builds a small corpus, captures a
mounted oracle, detaches the image, runs the proof parser, and compares output.

\subsection{What Not To Infer}

Do not infer native parser performance from the proof skeleton. The proof backend
currently shells out through \texttt{go run} and exists to keep correctness
experiments runnable while native root resolution and FS-record parsing are
split out.

\section{Native Rust Scanner}

\subsection{Current Boundary}

The native Rust crate is \path{crates/apfs-fastindex} and its
implementation-facing spec is \path{docs/implementation/rust-checkpoint-scanner.md}.
The crate is the v1 implementation slice; it covers the full read-only path
from source-gate to FS-tree record-family inventory but does not yet decode
FS-record bodies, emit \texttt{NamespaceEntry} rows, compute logical sizes, or
claim live-mounted-source correctness.

What the Rust path currently proves on the proof fixture:

\begin{itemize}
  \item Allowlisted input is a detached APFS \texttt{.dmg} image attached with
  \texttt{hdiutil attach -plist -nomount}, or a caller-supplied raw APFS
  container path under \path{/dev/disk*} or \path{/dev/rdisk*}.
  \item Block zero is validated as an APFS container locator (block size,
  \texttt{NXSB} magic, object type, object id, Fletcher-64 checksum).
  \item Contiguous checkpoint descriptor areas are scanned for checksum-valid
  \texttt{NXSB} candidates and the highest candidate XID is selected.
  \item The selected NX superblock is decoded end-to-end and the v1
  incompatible-feature allowlist is enforced.
  \item The per-checkpoint \texttt{checkpoint\_map\_phys\_t} ring is walked,
  ephemeral OID mappings are reported, and \texttt{CHECKPOINT\_MAP\_LAST} is
  required.
  \item The container OMAP is opened and queried with the
  \texttt{(oid, max\_xid)} lower-bound semantics required by SR-006, with
  hard stops on the encryption, no-header, and unknown-flag conditions.
  \item Every reachable volume superblock is decoded and gated by the v1
  allowlist (encryption, sealed volumes, normalization sensitivity, unknown
  incompatible features all force \texttt{Unsupported(reason)}).
  \item Supported volumes have their volume OMAP opened and the FS-tree root
  validated through it.
  \item The FS-tree is walked read-only and an FS-record family dump is
  emitted (\texttt{inode}, \texttt{xattr}, \texttt{sibling\_link},
  \texttt{dstream\_id}, \texttt{file\_extent}, \texttt{dir\_rec},
  \texttt{sibling\_map}, snapshot families, and an \texttt{unsupported}
  bucket).
\end{itemize}

\subsection{Module Layout}

\begin{verbatim}
crates/apfs-fastindex/src/
  lib.rs          orchestration, parser shapes, source gate
  main.rs         apfs-fastindex-scan binary (JSON output)
  block_io.rs     block reads, le_u*, Fletcher-64, test helpers
  object.rs       obj_phys_t validation (SR-007 single chokepoint)
  btree.rs        read-only B-tree node reader (fixed and variable kv)
  omap.rs         OMAP resolver: (oid, max_xid) lower-bound + summary
  container.rs    NXSB decode + checkpoint map ring walk
  volume.rs       APSB decode + v1 feature allowlist (SR-012)
  fs_records.rs   FS-tree record-family dumper (SR-008)
\end{verbatim}

The fail-closed contract is centralized in
\texttt{object::validate\_object\_block}: checksum, expected object type,
expected storage class, encryption flag, no-header flag, optional
\texttt{oid == paddr} check, and optional \texttt{xid <= scan\_xid} guard.
Every reader in the crate validates a block through this function before
trusting its body.

\subsection{Commands}

Run unit tests:

\begin{verbatim}
cargo test -p apfs-fastindex
\end{verbatim}

Build and run the scanner against a detached APFS \texttt{.dmg}:

\begin{verbatim}
cargo run --bin apfs-fastindex-scan -- /path/to/detached-apfs.dmg
\end{verbatim}

Run the EX-10 probe (also asserts on the native dump structure):

\begin{verbatim}
python3 docs/research/experiments/EX-10-rust-checkpoint-scanner/artifacts/probe_ex10.py
\end{verbatim}

Run the EX-12 OMAP lookup contract probe:

\begin{verbatim}
python3 docs/research/experiments/EX-12-omap-lookup-contract/artifacts/probe_ex12.py
\end{verbatim}

\subsection{Empirical Corrections Distilled Into Code}

Three empirical corrections from running the Rust path against the real proof
fixture have been landed in the crate and recorded in RL-01, RL-02, and RL-03:

\begin{enumerate}
  \item \texttt{OBJECT\_TYPE\_CHECKPOINT\_MAP} blocks carry \texttt{OBJ\_PHYSICAL},
  not \texttt{OBJ\_EPHEMERAL}, so the validator for the checkpoint-map ring
  expects \texttt{ExpectedStorage::Physical}.
  \item B-tree TOC \texttt{v\_off} measures the distance from the end of the
  value area to the start of the value; values run forward by
  \texttt{fixed\_value\_size} or \texttt{v\_len}. The earlier
  ``value extends backward from \texttt{v\_off}'' assumption mis-decoded OMAP
  values as \texttt{flags=0x10ffff}.
  \item The FS-tree root is reached through the volume OMAP, so on disk it
  carries a virtual \texttt{obj\_phys\_t} (storage flags 0). The FS-tree
  validator therefore uses \texttt{ExpectedStorage::Virtual} and checks
  \texttt{header.oid == volume\_omap\_lookup.oid} rather than requiring
  \texttt{o\_oid == paddr}.
\end{enumerate}

\section{Native Replacement Path}

Remaining native work, in order, with each step landing as its own
\texttt{EX-*} note before promotion:

\begin{enumerate}
  \item Resolve the \texttt{EX-14} block-1031 context blocker before any new
  body decoder is added.
  \item Replay \texttt{EX-13} with the SR-015 source-backed xfield cursor rule
  and \texttt{xf\_used\_data} validation.
  \item Add synthetic negative record-body cases from SR-016.
  \item Decode \texttt{j\_drec\_*\_key\_t} and \texttt{j\_drec\_val\_t} to
  reconstruct the path table without yet emitting \texttt{NamespaceEntry} rows.
  \item Decode \texttt{j\_inode\_val\_t} flags, link count, and embedded fields
  needed for logical-size selection per SR-009 and SR-017.
  \item Decode \texttt{j\_dstream\_id\_val\_t} and \texttt{j\_dstream\_t} for
  logical size, and resolve the symlink-target \texttt{XATTR}.
  \item Resolve \texttt{j\_sibling\_link\_val\_t} and
  \texttt{j\_sibling\_map\_val\_t} for hard-link unification.
  \item Wire the resolved namespace into the existing
  \texttt{ParserOutput} / \texttt{NamespaceEntry} / \texttt{DirectoryAggregate}
  / oracle-diff contract so the Rust path can be diffed against the Python
  proof oracle.
  \item Re-run the existing Python proof fixture and the expanded EX-04 corpus
  after each replacement stage. Native parity is not declared until the
  regression suite passes.
\end{enumerate}

Any support broadening must first update the related \texttt{SR-*},
\texttt{EX-*}, and \texttt{RL-*} artifacts.
